% Parked from 02-v2 main.tex on 2026-10-31 (author's decision: "middle" shortening).
% Contents: the typing-level results on typings onto U_1 (seed of a separate paper on S1/S2),
% the three-letter version of the linear-factor lower bound, the target separations of
% Appendix K (second half), and the rank-one appendix. All statements were proved and
% externally checked while in the manuscript (see README change log, reviews/).
% This file is not compiled; labels are kept unchanged so that blocks can be restored.

%==============================================================================
% Section 6: Proposition linear-factor-sharp (three letters, fan-out one) and surrounding text
%==============================================================================
The upper side of Corollary~\ref{cor:typed-sandwich} grows with $\|G\|$ and
the lower side does not.  At fan-out one the product $\|G\|\bar\theta_G$ is
nevertheless necessary, under every typing over an alphabet of at least three
letters and for presentations within a constant factor of the smallest.

\begin{proposition}[The linear factor is necessary at fan-out one]
\label{prop:linear-factor-sharp}
Let $\mathtt a,\mathtt c,\mathtt d\in\Sigma$ be distinct and let $s$ be any
typing.  For integers $q\ge k\ge1$ put $M:=2^k$, $W:=M+2q+2$ and
\[
 w_i:=\mathtt a^i\mathtt d\mathtt a^{q-i}\mathtt c^M\mathtt a^{q-i}\mathtt d
 \mathtt a^i\quad(0\le i\le q),\qquad L_{q,k}:=\{w_0,\ldots,w_q\}.
\]
\begin{enumerate}[label=\textup{(\roman*)},leftmargin=*]
\item Every subset of $L_{q,k}$ belongs to
$\Cmcf{1}{s_{\mathrm{triv}}}\subseteq\Cmcf{1}{s}$.
\item For $\mathcal B_{1,s}$, for $\mathcal B^{\le r}_{1,s}$ \textup($r\ge1$\textup),
and for every set-driven operator with characteristic samples throughout
$\Cmcf{1}{s}$, the only characteristic sample of $L_{q,k}$ is $L_{q,k}$ itself.
\item $L_{q,k}$ has a reduced good fan-out-one presentation $G_{q,k}$ of rule
rank two with $\|G_{q,k}\|=17q+6k+13\le23(q+1)$ and
$\bar\theta_{G_{q,k}}=t_{\mathrm{typ}}(G_{q,k},s)=W$.  Every reduced good
fan-out-one presentation $G$ of $L_{q,k}$ has $\bar\theta_G=W$ and
$\|G\|\ge q+1$.
\end{enumerate}
Consequently every characteristic sample $F$ in~\textup{(ii)} satisfies
\[
 \|F\|_+=(q+1)(W+1)\ \ge\ \tfrac1{23}\,\|G_{q,k}\|\,\bar\theta_{G_{q,k}}.
\]
\end{proposition}

The proof is in Appendix~\ref{app:linear-factor-fanout}.

With $q=k=n$ one has $\|G_{n,n}\|=23n+13$ and $\bar\theta_{G_{n,n}}=2^n+2n+2$;
with $k$ fixed, $\bar\theta_{G_{q,k}}=\Theta(\|G_{q,k}\|)$.  By (iii) the
presentations are within a constant factor of the smallest reduced good ones,
so the bound is not obtained by padding $G$.  With trivial typing the
proposition shows that Yoshinaka's bound $\|K_G\|\le|P|\theta_G$ is attained
up to a constant factor.  Since $t_{\mathrm{typ}}=\bar\theta_G$ in this family,
it says nothing about the further factor $\|G\|$ in
Theorems~\ref{thm:window-group-polydata} and~\ref{thm:local-group-polydata}.
The family is specific to fan-out one: in $L_{2,1}$ the pairs
$(\mathtt d,\mathtt{ac})$ and $(\mathtt{da},\mathtt c)$ share the context
$\blank_1\mathtt a\blank_2\mathtt{caad}$, but only the second is accepted by
$\blank_1\mathtt{ac}\blank_2\mathtt{aad}$, so $L_{2,1}$ is not
$(2,s_{\mathrm{triv}})$-tuple-substitutable.

%==============================================================================
% Section 6: paragraph after Proposition linear-factor-all-fanout
%==============================================================================
For $f=1$ this is a variant of Proposition~\ref{prop:linear-factor-sharp} over a
larger alphabet; for $f\ge2$ it removes the restriction to fan-out one.  In
$w_i$ the position of one marker $\mathtt d_h$, and the distance of two, determine
$i$.  A context shared by two different words carries at most $d-1$ of the $2f$
markers, so some component of each tuple carries two markers and identifies the
word.  As in Proposition~\ref{prop:linear-factor-sharp},
$t_{\mathrm{typ}}=\bar\theta_G=W$ for this family.  The proof is in
Appendix~\ref{app:linear-factor-fanout}.

%==============================================================================
% Section 7 tail: three-letter family, typing-level obstruction, binary typing quadratic family, constant weight (weight-slice, all fan-out), closing paragraph
%==============================================================================
\begin{proposition}[A three-letter doubling family]
\label{prop:three-letter-family}
Let $\mathtt a,\mathtt b,\mathtt c$ be distinct letters, put $K^{(0)}:=\{\mathtt a\}$
and $K^{(i)}:=\mathtt bK^{(i-1)}K^{(i-1)}\mathtt b\cup\mathtt cK^{(i-1)}\mathtt c$ for
$i\ge1$, and let $n\ge1$.  Then $K_n:=K^{(n)}$ is Clark--Eyraud substitutable, so
$K_n\in\Cmcf{1}{s}$ for every typing $s$ on every alphabet containing these
letters.  The grammar $G'_n$ with rules $S\to E_n$, $E_0\to\mathtt a$ and
$E_i\to\mathtt bE_{i-1}E_{i-1}\mathtt b\mid\mathtt cE_{i-1}\mathtt c$
($1\le i\le n$) is a reduced good rank-two presentation of $K_n$ with
$\|G'_n\|=13n+7$ and $\bar\theta_{G'_n}=4n$.  The only word of $K_n$ without
$\mathtt c$ is $V_n$, where $V_0:=\mathtt a$ and
$V_i:=\mathtt bV_{i-1}V_{i-1}\mathtt b$, and $|V_n|=3\cdot2^n-2$.
\end{proposition}

The proof is in Appendix~\ref{app:three-letter-family}.


\begin{corollary}[A typing-level obstruction at fan-out one]
\label{cor:typing-level-obstruction}
Let $r_{\max}\ge2$ and let $s:\Sigma^+\twoheadrightarrow\mathsf S$ be a typing.
Suppose that distinct idempotents $e,z\in\mathsf S$ satisfy $ez=ze=z$ and that
distinct letters $\mathtt a,\mathtt b,\mathtt c\in\Sigma$ satisfy
$s(\mathtt a)=s(\mathtt b)=e$ and $s(\mathtt c)=z$.  Then every characteristic
sample of $K_n$ for $\mathcal B^{\le r_{\max}}_{1,s}$, for $\mathcal B_{1,s}$, or for
any set-driven operator having characteristic samples throughout $\Cmcf{1}{s}$
has total encoded length at least $3\cdot2^n-1$, while $\|G'_n\|$ and
$\bar\theta_{G'_n}$ are $O(n)$.  So for this single typing no polynomial
bounds the characteristic-data length at fan-out one and rank $r_{\max}$ as in
Definition~\ref{def:uniform-polydata}.  This applies to every typing onto $U_1$
with at least two letters of type $1$.
\end{corollary}

\begin{proof}
In $\{e,z\}$, $e$ is an identity and $z$ is absorbing, so a word of $K_n$ has type
$e$ exactly when it contains no $\mathtt c$; by
Proposition~\ref{prop:three-letter-family} the only such word is $V_n$.
Proposition~\ref{prop:rare-type} forces $V_n$ into every such sample.  For a
typing onto $U_1$, surjectivity gives a letter of type $\zeta$, since words over
letters of type $1$ have type $1$; take $e=1$ and $z=\zeta$.
\end{proof}

Let $\Sigma=\{\mathtt a,\mathtt b\}$ and let $s_{U,2}:\Sigma^+\twoheadrightarrow U_1$ send $\mathtt a$ to
$1$ and $\mathtt b$ to $\zeta$.  For this typing, with a single letter of type $1$,
Corollary~\ref{cor:typing-level-obstruction} does not apply; still, its data can
be quadratic in $\bar\theta_G$ for presentations of logarithmic size.

\begin{proposition}[Quadratic data for the binary typing at fan-out one]
\label{prop:binary-u1-quadratic}
For $m,n\ge1$ let $L_{m,n}\subseteq\{\mathtt a,\mathtt b\}^+$ be the set of all
concatenations of $2^n$ words from $\{\mathtt a^{2^m},\mathtt b\}$.
\begin{enumerate}[label=\textup{(\roman*)},leftmargin=*]
\item $L_{m,n}\in\Cmcf{1}{s_{U,2}}$.
\item $L_{m,n}$ has a reduced good fan-out-one presentation $G_{m,n}$ of rule rank
two with $\|G_{m,n}\|=6m+6n+9$ and $\bar\theta_{G_{m,n}}=2^m+2^n-1$.
\item For $\mathcal B_{1,s_{U,2}}$, for $\mathcal B^{\le r}_{1,s_{U,2}}$ \textup($r\ge1$\textup),
and for every set-driven operator with characteristic samples throughout
$\Cmcf{1}{s_{U,2}}$, every characteristic sample $F$ of $L_{m,n}$ has
$\|F\|_+>2^{m+n}$.
\end{enumerate}
Consequently, for $m=n$ one has $\|G_{n,n}\|=12n+9$, $\bar\theta_{G_{n,n}}=2^{n+1}-1$
and $\|F\|_+>\bar\theta_{G_{n,n}}^{\,2}/4$; in particular, for this typing at
fan-out one there is no bound of the form $p(\|G\|)\,\bar\theta_G$ with $p$ a
polynomial, in contrast with Theorem~\ref{thm:local-group-polydata}.
\end{proposition}

The proof is in Appendix~\ref{app:binary-u1-quadratic}.  The pure word
$\mathtt a^{2^{m+n}}$ is the product of the block length and the number of
blocks, while $\bar\theta_G$ is about their sum, since a shortest sentence using
the long block fills every other position with $\mathtt b$.  The family is not
$(2,s_{U,2})$-tuple-substitutable for $m\ge3$ and $n\ge2$, since it has the factor
$\mathtt a^{2^m}\mathtt b\mathtt a^{2^m}$ but is not Parikh-saturated
(Lemma~\ref{lem:app-u1-deep}); so it does not contradict
Theorem~\ref{thm:binary-u1-polydata}.  Whether the data for this typing at
fan-out one are polynomial remains open.

The exponent two is sharp on a natural subclass that contains these targets.
For $w\in\{\mathtt a,\mathtt b\}^*$ and $\lambda>0$ put
$\psi_\lambda(w):=|w|_{\mathtt a}+\lambda|w|_{\mathtt b}$; a language has
\emph{constant weight} if $\psi_\lambda$ is constant on it for some $\lambda>0$.
The languages $L_{m,n}$ have constant weight with $\lambda=2^m$.

\begin{proposition}[Sharp quadratic bound for constant weight]
\label{prop:binary-u1-weight-slice}
Let $G$ be a reduced good fan-out-one grammar over $\{\mathtt a,\mathtt b\}$ with
$v$ nonterminals and $\theta:=\bar\theta_G$, and suppose that $L(G)$ has constant
weight.  Put $B(\theta):=\lfloor(\theta+1)^2/4\rfloor$.
\begin{enumerate}[label=\textup{(\roman*)},leftmargin=*]
\item Every nonterminal with a pure yield has exactly one pure yield, of
length at most $B(\theta)$; hence $\mu^\star(L(G),s_{U,2})\le B(\theta)$.
\item $t_{\mathrm{typ}}(G,s_{U,2})\le\max\{B(\theta),v(\theta-1)+1\}$.
\item If moreover $L\in\Cmcf{1}{s_{U,2}}$, $L(G)=L\cap\Sigma^+$ and $G$ has rule rank
at most $r$, then $\mathcal B^{\le r}_{1,s_{U,2}}$ and $\mathcal B_{1,s_{U,2}}$ have a
common characteristic sample $K$ with
$\|K\|_+\le3^r(r+3)\|G\|^2\theta^2+\mathbf 1_{\{\eps\in L\}}$.
\end{enumerate}
Items \textup{(i)} and \textup{(ii)} do not use substitutability.  For the
presentations $G_{n,n}$ of Proposition~\ref{prop:binary-u1-quadratic} the bound in
\textup{(i)} is attained: $\mu^\star=4^n=(\bar\theta_{G_{n,n}}+1)^2/4$.
\end{proposition}

The proof is in Appendix~\ref{app:binary-u1-quadratic}.  The key step for (i):
under constant weight $\psi_\lambda\equiv W$, a pure yield is determined by its
weight, and a lowest node of pure yield length $h\ge\lambda$ has a fixed yield, so
every sentence through its rule has length at least $h+(W-h)/\lambda$; hence
$W\le\lambda(\bar\theta_G+1-\lambda)\le(\bar\theta_G+1)^2/4$.


The bound on the shortest word of type $1$ does not depend on the fan-out or on
the number of letters.  For a typing $s$ onto $U_1$ put
$\Sigma_1:=\{x\in\Sigma:s(x)=1\}$, and for $\lambda>0$ let $\psi_\lambda(w)$ be
the number of letters of $w$ in $\Sigma_1$ plus $\lambda$ times the number of the
other letters.

\begin{proposition}[Constant weight at every fan-out]
\label{prop:u1-weight-all-fanout}
Let $s$ be a typing onto $U_1$ and let $G$ be a reduced good MCFG, of any fan-out
and rank, such that $\psi_\lambda$ is constant on $L(G)$ for some $\lambda>0$.
Then $\mu^\star(L(G),s)\le\lfloor(\bar\theta_G+1)^2/4\rfloor$.
\end{proposition}

Substitutability is not used.  For $s_{U,2}$ at fan-out one this is the bound on
$\mu^\star$ in Proposition~\ref{prop:binary-u1-weight-slice}(i), attained by
$G_{n,n}$.  The proof, in Appendix~\ref{app:binary-u1-quadratic}, does not use
uniqueness of pure yields; it compares a shortest derivation of a word of type
$1$ with shortest successful uses of its rules, which works for tuples of any
arity.

Constant weight is not necessary for $(1,s_{U,2})$-tuple-substitutability: the
language $\{\mathtt{aa},\mathtt{ab},\mathtt{bbb}\}$ is substitutable but has no
constant weight.  For typings onto $U_1$ at fan-out one the picture is
therefore: two letters of type $1$ force exponential data
(Corollary~\ref{cor:typing-level-obstruction}); a single letter of type $1$
forces data quadratic in $\bar\theta_G$, which is sharp for constant weight
(Propositions~\ref{prop:binary-u1-quadratic}
and~\ref{prop:binary-u1-weight-slice}), and whether it admits polynomial data
in general is Open Problem~\ref{op:S1}.

%==============================================================================
% Conclusion: old Open Problems S1 and fanout-two
%==============================================================================
\begin{openproblem}[The single-neutral-letter typing at fan-out one]
\label{op:S1}
For $s_{U,2}$, fan-out one and $r_{\max}\ge2$: is the characteristic-data length
bounded by a polynomial in $\|G\|$ and $\bar\theta_G$?  By
Proposition~\ref{prop:binary-u1-quadratic} such a polynomial must have degree at
least two in $\bar\theta_G$, and by Proposition~\ref{prop:binary-u1-weight-slice}
degree two suffices for targets of constant weight.  Two words
$\mathtt a^p,\mathtt a^q$ ($p<q$) of a target share the context $(\eps,\eps)$ and
force $\mathtt a^{p+k(q-p)}\in L$ for all $k\ge0$, so a finite target has at most
one word of type $1$.  A negative answer, together
with Theorem~\ref{thm:binary-u1-polydata}, would show that one typing needs
superpolynomial data at fan-out one and polynomial data at fan-out two.
\end{openproblem}

\begin{openproblem}[Pseudovarieties at higher fan-out]
\label{op:fanout-two}
Fix $f\ge2$ and $r_{\max}\ge2$.  For which typings $s$ with
$\mathsf S\notin\mathbf{LG}$ is there a polynomial $p$ such that, for every
$L\in\Cmcf{f}{s}$ and every reduced good fan-out-$\le f$, rank-$\le r_{\max}$
presentation $G$ of $L\cap\Sigma^+$, the operator
$\mathcal B^{\le r_{\max}}_{f,s}$ has a characteristic sample of total length
at most $p(\|G\|,\bar\theta_G)$?  At $f=2$ the typing $s_{U,2}$ is such a typing
(Theorem~\ref{thm:binary-u1-polydata}); typings onto $U_1$ with two letters of
type $1$, which need exponential data at fan-out one
(Corollary~\ref{cor:typing-level-obstruction}), and all $f\ge3$ are open.  Every
pseudovariety not contained in $\mathbf{LG}$ contains $U_1$, so a
superpolynomial lower bound for one typing onto $U_1$ at fan-out $f$ would make
$\mathbf{LG}$ the boundary at that fan-out as well.  By
Proposition~\ref{prop:rare-type}, a family with $\mu^\star(L_n,s)=n^{\omega(1)}$
and $\|G_n\|\bar\theta_{G_n}=n^{O(1)}$ would suffice;
Theorem~\ref{thm:doubling-no-lift} excludes the bracketed doubling family, and
Proposition~\ref{prop:u1-weight-all-fanout} excludes every family of constant
weight.  Conversely, polynomial data at $f=2$ for a typing onto $U_1$ with two
letters of type $1$ would show that one typing needs exponential data at fan-out
one and polynomial data at fan-out two.
\end{openproblem}

%==============================================================================
% Introduction: old Theorem B with items (ii), (iii)
%==============================================================================
\begin{mainthm}[The boundary at fan-out one]
\label{thm:B}
Let $f=1$.
\begin{enumerate}[label=\textup{(\roman*)},leftmargin=*]
\item A pseudovariety of finite semigroups admits uniform polynomial
characteristic-data length iff it is contained in $\mathbf{LG}$
\textup(Theorem~\ref{thm:fanout-one-data}\textup).
\item If some copy of $U_1$ in $\mathsf S$ is realized by two letters of its
identity type and one letter of its absorbing type, the class needs
exponential characteristic data
\textup(Corollary~\ref{cor:typing-level-obstruction}\textup).
\item For $\Sigma=\{\mathtt a,\mathtt b\}$ and the typing $s_{U,2}$ with
$\mathtt a\mapsto1$, $\mathtt b\mapsto\zeta$ onto $U_1$, there are targets with
$\|G\|=O(\log\bar\theta_G)$ whose characteristic samples exceed
$\bar\theta_G^{\,2}/4$, so no bound $p(\|G\|)\,\bar\theta_G$ with $p$ a polynomial
holds \textup(Proposition~\ref{prop:binary-u1-quadratic}\textup).  For targets
whose words all have the same positive weight $|w|_{\mathtt a}+\lambda|w|_{\mathtt b}$,
every realized type has a word of length at most $(\bar\theta_G+1)^2/4$, with
equality for the targets of the first sentence \textup(Proposition~\ref{prop:binary-u1-weight-slice}\textup).
More generally, for every typing onto $U_1$ and every fan-out, the same holds
whenever letters of type $1$ have weight $1$, all other letters a common weight
$\lambda>0$, and all words of the target the same weight
\textup(Proposition~\ref{prop:u1-weight-all-fanout}\textup).
\end{enumerate}
\end{mainthm}

%==============================================================================
% Introduction: old remaining-question paragraph
%==============================================================================
\paragraph{The remaining question.}
Theorems~\ref{thm:B}(iii) and~\ref{thm:C}(i) bracket one typing: at fan-out one
it needs data quadratic in $\bar\theta_G$, at fan-out two only polynomial data.
Whether its data at fan-out one are polynomial is the central open problem
(Open Problem~\ref{op:S1}).  A negative answer would show that raising the
fan-out can make learning strictly cheaper; a positive answer would classify
all typings onto $U_1$ at fan-out one by their number of letters of type $1$.
The same separation would follow from polynomial data at fan-out two for a
typing onto $U_1$ with two letters of type $1$, since such typings need
exponential data at fan-out one (Open Problem~\ref{op:fanout-two}).  Thus
Theorem~\ref{thm:C} shows that the fan-out-one obstruction does not lift, but
not yet that the boundary itself moves.

%==============================================================================
% Appendix I: proof of Proposition linear-factor-sharp
%==============================================================================
\begin{proof}[Proof of Proposition~\ref{prop:linear-factor-sharp}]
(i) Let $J\subseteq L_{q,k}$ and $x,y\in\Sigma^+$ with $uxv,uyv\in J$.  For
trivial typing and $d=f=1$ we must show $\D_J^{(1)}(x)=\D_J^{(1)}(y)$.  All
words of $L_{q,k}$ have length $W$, so $|x|=|y|$.  The two occurrences of
$\mathtt d$ in $w_i$ are at positions $i+1\in[1,q+1]$ and
$W-i\in[W-q,W]$; these ranges are disjoint since $W-q=M+q+2$, so the
position of a single occurrence of $\mathtt d$ determines $i$.  If $u$ or $v$
contains $\mathtt d$, that occurrence has the same position in $uxv$ and in
$uyv$, hence $uxv=uyv$ and $x=y$.  Otherwise both occurrences of
$\mathtt d$ lie in $x$; their distance $W-2i-1$ determines $i$, and the
offset of the first $\mathtt d$ in $x$ determines the position of $x$.  So
$x$ has exactly one occurrence in the words of $L_{q,k}$, and
$\D_J^{(1)}(x)=\{(u,v)\}$; likewise $\D_J^{(1)}(y)=\{(u,v)\}$.  Finite
languages are context-free, and the inclusion is item~(a) after
Definition~\ref{def:class-cmcf}.

(ii) Let $F$ be a characteristic sample with $w_i\notin F$.  Then
$F\subseteq L':=L_{q,k}\setminus\{w_i\}$, and $L'\in\Cmcf{1}{s}$ by (i).
Soundness (Proposition~\ref{prop:reconstruction-soundness}) for the target
$L'$ gives $L(\mathcal B(F))\subseteq L'\subsetneq L_{q,k}$ for both canonical
operators.  For a class-wide set-driven operator, a characteristic sample
$F'$ of $L'$ gives a finite $F\cup F'\subseteq L'$ on which the output must
equal both $L_{q,k}$ and $L'$, as in the proof of
Proposition~\ref{prop:rare-type}.  Hence $F=L_{q,k}$.  Conversely,
$K=L_{q,k}$ is the only finite $K$ with $L_{q,k}\subseteq K\subseteq L_{q,k}$,
and the canonical operators return $L_{q,k}$ on it by sample consistency and
soundness.

(iii) Take nonterminals $A_0,\ldots,A_q$, $H_0,\ldots,H_q$, $C_0,\ldots,C_k$,
start symbol $A_q$ and the rules
\[
\begin{array}{lll}
 A_j\to\mathtt d H_j\mathtt d\ (0\le j\le q), &
 A_j\to\mathtt a A_{j-1}\mathtt a\ (1\le j\le q), &
 H_0\to C_k,\\
 H_j\to\mathtt a H_{j-1}\mathtt a\ (1\le j\le q), &
 C_0\to\mathtt c, &
 C_t\to C_{t-1}C_{t-1}\ (1\le t\le k).
\end{array}
\]
By induction, $L(C_t)=\{\mathtt c^{2^t}\}$,
$L(H_j)=\{\mathtt a^j\mathtt c^M\mathtt a^j\}$ and
$L(A_j)=\{\mathtt a^i\mathtt d\mathtt a^{j-i}\mathtt c^M\mathtt a^{j-i}
\mathtt d\mathtt a^i:0\le i\le j\}$, so $L(G_{q,k})=L_{q,k}$; the grammar is
good and reduced.  It has $2q+k+3$ nonterminals and, in the order listed,
$q+1,q,1,q,1,k$ rules of sizes $5,5,3,5,2,5$
(Definition~\ref{def:standard-mcfg}), which gives $\|G_{q,k}\|=17q+6k+13$;
$k\le q$ gives the bound $23(q+1)$.  In a reduced presentation every rule
occurs in a successful derivation, and all words of $L_{q,k}$ have length
$W$, so $\bar\theta_G=W$.  In $G_{q,k}$ all yields of a nonterminal, and hence
of each typed nonterminal, have one length, at most $W$; the typed start
nonterminals have yields of length $W$, so $t_{\mathrm{typ}}(G_{q,k},s)=W$.
Finally let $G$ be any reduced good fan-out-one presentation of $L_{q,k}$.
For trivial typing, Theorem~\ref{thm:exact-reconstruction} makes
$\Wit(\widetilde G_0)$ characteristic for $\mathcal B_{1,s_{\mathrm{triv}}}$,
and it has at most $|V|+|P|\le\|G\|$ words
(Definition~\ref{def:witness-set} and the remark after
Theorem~\ref{thm:data-from-typed-thickness}); by (ii) for
$s_{\mathrm{triv}}$, it has $q+1$ words.

The final estimate follows from (ii), (iii) and $|w_i|=W$.
\end{proof}

%==============================================================================
% Appendix I: Lemma app-lf-descent (used only by the parked quadratic family)
%==============================================================================
The following lemma is used in Appendix~\ref{app:binary-u1-quadratic}.

\begin{lemma}\label{lem:app-lf-descent}
Let $G$ be a reduced good MCFG with $L(G)\ne\emptyset$, and let $\mu$ be the least
length of a word of $L(G)$.  For every $w\in L(G)$ with $|w|>\mu$ there is
$w'\in L(G)$ with $1\le|w|-|w'|\le\bar\theta_G-\mu$.
\end{lemma}

\begin{proof}
For $B\in V$ let $d_B$ be the least total length of a tuple in $L_B(G)$ and $o_B$
the least number of terminals outside one $B$-subtree of a successful
derivation.  By linearity and nondeletion, a rule $R:B\to\rho(B_1,\ldots,B_r)$
writing $c_R$ terminals has $\theta_G(R)=o_B+c_R+\sum_jd_{B_j}$, since the outside,
the rule and the children can be chosen independently; moreover
$o_B+d_B\ge\mu$, since replacing the $B$-subtree of a derivation that attains $o_B$
by a derivation of a shortest $B$-tuple gives a word of length $o_B+d_B$.  In a derivation of $w$ choose a deepest node $u$ whose yield is
longer than $d_{B_u}$; the root is such a node.  Its children have yields of
length $d_{B_j}$, so replacing the subtree at $u$ by a derivation of a shortest
$B_u$-tuple shortens $w$ by
$c_R+\sum_jd_{B_j}-d_{B_u}=\theta_G(R)-o_{B_u}-d_{B_u}\in[1,\bar\theta_G-\mu]$.
\end{proof}

%==============================================================================
% Appendix G: substitutability of the three-letter doubling family
%==============================================================================
\section{Substitutability of the Three-Letter Doubling Family}
\label{app:three-letter-family}

This appendix proves that $K_n$ in Proposition~\ref{prop:three-letter-family} is
Clark--Eyraud substitutable.  Fix $n\ge1$ and write $x\equiv y$ if
$\{(u,v):uxv\in K_n\}=\{(u,v):uyv\in K_n\}$, a monoid congruence.  Stacks are
words over $E_0,\ldots,E_n$ and the closing markers $\hat{\mathtt b}_i,\hat{\mathtt c}_i$
($1\le i\le n$), top on the left, and for every stack $\mathcal Z$ and $1\le i\le n$
the moves are
\[
\begin{aligned}
E_0\mathcal Z&\xrightarrow{\mathtt a}\mathcal Z,&
E_i\mathcal Z&\xrightarrow{\mathtt b}E_{i-1}E_{i-1}\hat{\mathtt b}_i\mathcal Z,&
E_i\mathcal Z&\xrightarrow{\mathtt c}E_{i-1}\hat{\mathtt c}_i\mathcal Z,\\
\hat{\mathtt b}_i\mathcal Z&\xrightarrow{\mathtt b}\mathcal Z,&
\hat{\mathtt c}_i\mathcal Z&\xrightarrow{\mathtt c}\mathcal Z
\end{aligned}
\]
and no others.  The moves are deterministic, and properties (P2) and (P3) of
Appendix~\ref{app:bracketed-family} hold with the same proofs.  So does the first
part of Lemma~\ref{lem:app-br-recognition}: $E_i\mathcal Z\xrightarrow{w}\mathcal Z$ iff
$w\in K^{(i)}$.  Property (P1) fails, because $\mathtt b$ and $\mathtt c$ both open
and close: for instance $\mathtt b$ maps $E_i\mathcal Z$ and
$\hat{\mathtt b}_jE_{i-1}E_{i-1}\hat{\mathtt b}_i\mathcal Z$ to the same stack.  We
replace it by injectivity on reachable stacks, those $\mathcal X$ with
$E_n\xrightarrow{u}\mathcal X$ for some $u$.

\begin{lemma}\label{lem:app-tl-reachable}
Put $\mathcal F_i:=\{E_{i-1}\hat{\mathtt b}_i,\hat{\mathtt b}_i,\hat{\mathtt c}_i\}$.
Every reachable stack is $\eps$ or $\mathcal YF_{j+1}\cdots F_n$ with
$F_i\in\mathcal F_i$ and $\mathcal Y\in\{E_j,\hat{\mathtt b}_j,\hat{\mathtt c}_j\}$
($\mathcal Y=E_j$ with $0\le j\le n$, or a marker with $1\le j\le n$).  Each
letter, and hence each word, acts as a partial injection on reachable stacks.
\end{lemma}

\begin{proof}
The set of such stacks contains $E_n$ and is closed under the moves: from
$E_jF_{j+1}\cdots F_n$, $\mathtt b$ gives $E_{j-1}(E_{j-1}\hat{\mathtt b}_j)F_{j+1}\cdots F_n$
and $\mathtt c$ gives $E_{j-1}\hat{\mathtt c}_jF_{j+1}\cdots F_n$; popping the top
leaves $F_{j+1}\cdots F_n$, which is $\eps$ or again of the stated form, a first
frame $E_j\hat{\mathtt b}_{j+1}$ being read as $\mathcal Y=E_j$ followed by the frame
$\hat{\mathtt b}_{j+1}$.  For injectivity, $\mathtt a$ only pops $E_0$.  A
$\mathtt b$-opening produces a stack beginning with two symbols $E$, from which
$i$ and $\mathcal Z$ are recovered; a $\mathtt b$-closing produces a product of
frames, which begins with a marker or with $E_j\hat{\mathtt b}_{j+1}$, never with two
symbols $E$, and whose first marker (or emptiness, for $i=n$) determines the
popped level.  The same holds for $\mathtt c$, whose openings begin with
$E_{i-1}\hat{\mathtt c}_i$.
\end{proof}

\begin{lemma}\label{lem:app-tl-recognition}
If $\mathcal X$ is reachable, $w\in K^{(i)}$ and $\mathcal X\xrightarrow{w}\mathcal Y$, then
$\mathcal X=E_i\mathcal Y$.  Consequently $w\equiv w'$ for all $w,w'\in K^{(i)}$.
\end{lemma}

\begin{proof}
Induct on $i$; for $i=0$, $\mathtt a$ only pops $E_0$.  Let $i\ge1$; all stacks
met are reachable.  If $w=\mathtt cv\mathtt c$ with $v\in K^{(i-1)}$ and the first
$\mathtt c$ opens $E_j$, induction applied to $v$ at $E_{j-1}\hat{\mathtt c}_j\mathcal Z$
gives $j=i$, and the last $\mathtt c$ pops $\hat{\mathtt c}_i$.  If the first
$\mathtt c$ pops $\hat{\mathtt c}_j$, the rest is a product of frames; to read $v$ it
must begin with $E_{i-1}$, so its first frame is $E_{i-1}\hat{\mathtt b}_i$, and after
$v$ the top is $\hat{\mathtt b}_i$, on which $\mathtt c$ has no move.  If
$w=\mathtt bv_1v_2\mathtt b$ and the first $\mathtt b$ opens $E_j$, induction gives
$j=i$ and the two child words pop the two copies of $E_{i-1}$ before the last
$\mathtt b$ pops $\hat{\mathtt b}_i$.  If the first $\mathtt b$ pops $\hat{\mathtt b}_j$,
then as before $v_1$ leaves the top $\hat{\mathtt b}_i$, from which $v_2$ cannot be
read.  This proves the first claim.  For the second, if $uwv\in K_n$, the run
$E_n\xrightarrow{u}\mathcal X\xrightarrow{w}\mathcal Y\xrightarrow{v}\eps$ has
$\mathcal X=E_i\mathcal Y$, and $E_i\mathcal Y\xrightarrow{w'}\mathcal Y$; exchange $w,w'$
for the converse.
\end{proof}

\begin{proposition}\label{prop:app-tl-congruence}
$K_n$ is Clark--Eyraud substitutable.
\end{proposition}

\begin{proof}
Lemma~\ref{lem:app-br-branch} holds with the same proof, with $\mathtt c,\mathtt b$
in place of $\lceil,\lfloor$ and $K^{(i)}$ in place of $L^{(i)}$, since it uses only
determinism, (P3) and the recognition of complete words from $E_i$.  The proof
of the first claim of Proposition~\ref{prop:app-br-congruence} then gives: if
$\mathcal X\xrightarrow{x}\mathcal Y$ and $\mathcal X\xrightarrow{y}\mathcal Y$, then
$x\equiv y$, with Lemma~\ref{lem:app-tl-recognition} in place of
Lemma~\ref{lem:app-br-exchange}.  Now let $x,y$ be nonempty with
$uxv,uyv\in K_n$.  Determinism gives runs
$E_n\xrightarrow{u}\mathcal X\xrightarrow{x}\mathcal Y\xrightarrow{v}\eps$ and
$E_n\xrightarrow{u}\mathcal X\xrightarrow{y}\mathcal Y'\xrightarrow{v}\eps$; $\mathcal Y$ and
$\mathcal Y'$ are reachable, so Lemma~\ref{lem:app-tl-reachable} gives
$\mathcal Y=\mathcal Y'$, and $x\equiv y$.
\end{proof}

\begin{proof}[Proof of Proposition~\ref{prop:three-letter-family}]
Substitutability is proved in Appendix~\ref{app:three-letter-family}; letters
outside $\{\mathtt a,\mathtt b,\mathtt c\}$ only add empty distributions, and
finite languages are context-free.  The grammar is good, and reduced since
$E_i\Rightarrow^*\mathtt c^i\mathtt a\mathtt c^i$.  It has $n+2$ nonterminals
and rules of sizes $3$ (start), $2$ (base), $7$ and $5$ (the two rules of each
level), whence $\|G'_n\|=13n+7$.  A shortest yield of $E_i$ has length $2i+1$,
and the levels above $E_i$ add at least $2(n-i)$ letters, attained by unary
rules.  So the start, base and unary rules have rule thickness $2n+1$, and the
binary rule of level $i$ has $2(n-i)+2+2(2i-1)=2n+2i$; hence
$\bar\theta_{G'_n}=4n$.  A word without $\mathtt c$ uses only binary rules,
which gives $V_n$.
\end{proof}

%==============================================================================
% Appendix J: quadratic data for the binary typing at fan-out one; constant-weight proofs
%==============================================================================
\section{Quadratic Data for the Binary Typing at Fan-Out One}
\label{app:binary-u1-quadratic}

This appendix proves Proposition~\ref{prop:binary-u1-quadratic}.  Put
$k:=2^m$, $M:=2^n$ and $\psi(w):=|w|_{\mathtt b}+|w|_{\mathtt a}/k$, which is
additive.  For a word $w$ let $\lambda(w)$ and $\varrho(w)$ be the lengths of its
longest suffix and longest prefix in $\mathtt a^*$.

\begin{proof}[Proof of Proposition~\ref{prop:binary-u1-quadratic}]
(i) A word of $\Sigma^+$ lies in $L_{m,n}$ iff all its maximal runs of $\mathtt a$
have length divisible by $k$ and $\psi(w)=M$: such a word factors uniquely into
blocks $\mathtt a^k$ and $\mathtt b$, and $\psi$ counts the blocks.  Let
$x,y\in\Sigma^+$ have the same type and let $uxv,uyv\in L_{m,n}$.  By additivity
$\psi(x)=\psi(y)$.  If $x,y$ are pure, this gives $x=y$.  Otherwise write
$x=\mathtt a^{\alpha}x'\mathtt a^{\beta}$ with $x'$ beginning and ending with
$\mathtt b$, and similarly $y=\mathtt a^{\alpha'}y'\mathtt a^{\beta'}$.  The maximal
$\mathtt a$-runs inside $x'$ and $y'$ are maximal runs of $uxv$ and $uyv$, so their
lengths are divisible by $k$; the runs through the two ends of $x$ have lengths
$\lambda(u)+\alpha$ and $\beta+\varrho(v)$, so $\alpha\equiv\alpha'$ and
$\beta\equiv\beta'\pmod k$.  For any context $(p,q)$ the maximal $\mathtt a$-runs of
$pxq$ are those of $p$ not at its end, a run of length $\lambda(p)+\alpha$, the
runs inside $x'$, a run of length $\beta+\varrho(q)$, and those of $q$ not at its
start; the same holds for $pyq$ with $\alpha',\beta',y'$.  So the divisibility
conditions agree, $\psi(pxq)=\psi(pyq)$, and $pxq\in L_{m,n}\iff pyq\in L_{m,n}$.
Finite languages are context-free.

(ii) Take nonterminals $C_0,\ldots,C_m$, $A_0,\ldots,A_n$, start symbol $A_n$, and the
rules $C_0\to\mathtt a$, $C_j\to C_{j-1}C_{j-1}$ ($1\le j\le m$),
$A_0\to C_m\mid\mathtt b$ and $A_i\to A_{i-1}A_{i-1}$ ($1\le i\le n$).  Then
$L(C_j)=\{\mathtt a^{2^j}\}$, $L(A_0)=\{\mathtt a^k,\mathtt b\}$ and $L(A_i)$ is the
set of concatenations of $2^i$ words of $L(A_0)$, so $L(A_n)=L_{m,n}$; the grammar
is good and reduced.  It has $m+n+2$ nonterminals and rules of sizes $2$, $5$
($m$ rules), $3$, $2$ and $5$ ($n$ rules) (Definition~\ref{def:standard-mcfg}),
which gives $6m+6n+9$.  With $d_B$ and $o_B$ as in the proof of
Lemma~\ref{lem:app-lf-descent}, $\theta_G(R)=o_B+c_R+\sum_jd_{B_j}$, where
$d_{A_0}=1$, $d_{A_i}=2^i$, $d_{C_j}=2^j$, $o_{A_i}=M-2^i$, $o_{A_0}=o_{C_m}=M-1$ and
$o_{C_j}=M-1+k-2^j$.  This gives $\theta_G(R)=M$ for the rules of $A_i$ and for
$A_0\to\mathtt b$, and $\theta_G(R)=M+k-1$ for $A_0\to C_m$ and the rules of the
$C_j$; so $\bar\theta_{G_{m,n}}=2^m+2^n-1$.

(iii) The only pure word of $L_{m,n}$ is $\mathtt a^{kM}$, so by
Proposition~\ref{prop:rare-type} every characteristic sample contains it and
$\|F\|_+\ge kM+1>2^{m+n}$.  For $m=n$, $\bar\theta_{G_{n,n}}<2^{n+1}$ gives
$\bar\theta_{G_{n,n}}^{\,2}/4<4^n<\|F\|_+$.  A bound $p(\|G\|)\bar\theta_G$ would give
$4^n\le p(12n+9)\,2^{n+1}$ for all $n$, which fails for every polynomial $p$.
\end{proof}

\begin{proof}[Proof of Proposition~\ref{prop:binary-u1-weight-slice}]
Let $\psi:=\psi_\lambda$ be constant, equal to $W>0$, on $L(G)$.  A fixed
successful context of a nonterminal $A$ shows that all yields of $A$ have one
weight $w_A$, so a pure yield of $A$ is the unique word $\mathtt a^{\ell_A}$ with
$\ell_A=w_A\le W$.  By Lemma~\ref{lem:thickness-equivalence} the shortest yield
of $A$ has length at most $\theta$.

(i) If $\lambda\le1$, the weight of that shortest yield gives $\ell_A\le\theta$.
If $\lambda>1$ and $\ell_A<\lambda$, no yield of $A$ contains $\mathtt b$, so again
$\ell_A\le\theta$.  Let $\lambda>1$ and $\ell_A\ge\lambda$.  In a derivation of
$\mathtt a^{\ell_A}$ from $A$ choose a node whose yield has length $h\ge\lambda$ and
whose children have yields of length less than $\lambda$; one is found by
starting at the root and moving to a child of yield length at least $\lambda$ as
long as there is one.  Let $R$ be its rule.  A child $C$ has weight $w_C<\lambda$, so it has no yield containing
$\mathtt b$ and its only yield is pure; $R$ writes no $\mathtt b$.  Hence every
application of $R$ yields the same pure word of length $h$.  In a shortest
sentence using $R$ the outside has weight $W-h$ and therefore length at least
$(W-h)/\lambda$, since no letter weighs more than $\lambda$.  So
\[
 \theta\ge h+\frac{W-h}{\lambda}\ge\frac W\lambda+\lambda-1,
 \qquad
 \ell_A\le W\le\lambda(\theta+1-\lambda)\le\frac{(\theta+1)^2}{4}.
\]
Since $\theta\le B(\theta)$, all cases give $\ell_A\le B(\theta)$.  For $A=S$
this bounds $\mu_1$.  If some word of $L(G)$ contains $\mathtt b$, some rule writes
$\mathtt b$, and every sentence using it has type $\zeta$, so
$\mu_\zeta\le\theta\le B(\theta)$.  Hence $\mu^\star(L(G),s_{U,2})\le B(\theta)$.

(ii) Pure typed nonterminals are covered by (i).  If $A$ has a yield containing
$\mathtt b$, choose a shortest path in the dependency graph from $A$ to the
left-hand side of a rule writing $\mathtt b$, and complete every off-path child by
a shortest yield.  By the formula
$\theta_G(R)=o_B+c_R+\sum_jd_{B_j}$ from the proof of
Lemma~\ref{lem:app-lf-descent}, the last rule contributes at most $\theta$
letters and each earlier node at most $\theta-d_C\le\theta-1$, where $C$ is its
path child; the path has at most $v$ nodes.

(iii) By (i), (ii) and $\tau\le t_{\mathrm{typ}}$,
Corollary~\ref{cor:distribution-representatives} gives
$\|K\|_+\le N_{\mathrm{typ}}(1+2\theta+r\max\{B(\theta),v\theta\})+\mathbf 1_{\{\eps\in L\}}$
with $N_{\mathrm{typ}}\le3^r\|G\|$; use $B(\theta)\le\theta^2$, $v\le\|G\|$ and
$1+2\theta\le3\theta^2$.  For $G_{n,n}$, $\bar\theta=2^{n+1}-1$ and $\mu^\star=4^n$ by
Proposition~\ref{prop:binary-u1-quadratic}.
\end{proof}

\begin{proof}[Proof of Proposition~\ref{prop:u1-weight-all-fanout}]
Put $\theta:=\bar\theta_G$, $\psi:=\psi_\lambda$, extended additively to tuples and
to the fixed parts of contexts, and let $W$ be its value on $L(G)$.  Some word of
$L(G)$ has type $\zeta$ only if some rule writes a letter outside $\Sigma_1$;
every sentence using that rule has type $\zeta$, so $\mu_\zeta\le\theta$.  Let
$w$ be a shortest word of type $1$ in $L(G)$, if there is one, and put
$N:=|w|=\psi(w)=W$.  Fix a successful derivation $T$ of $w$, and let $\ell_u$ be
the total length of the tuple derived at a node $u$ of $T$; all these tuples have
only letters in $\Sigma_1$.  For each rule $R$ fix a successful derivation of a
word $Z_R$ of length $\theta_G(R)\le\theta$ that uses $R$, an occurrence of $R$
in it, the tuple $\boldsymbol\eta_R$ derived there, the tuples
$\boldsymbol\eta_{R,j}$ of its children, and the context $D_R$ obtained by
deleting that subtree.  Let $u$ be a node of $T$ that uses $R$, with children
$u_j$.  Replacing the subtree at $u_j$ by a derivation of $\boldsymbol\eta_{R,j}$
gives a successful derivation; by linearity and nondeletion its word has weight
$N-\ell_{u_j}+\psi(\boldsymbol\eta_{R,j})$, so
$\psi(\boldsymbol\eta_{R,j})=\ell_{u_j}$.  Replacing the designated subtree of
$Z_R$ by the subtree at $u$ gives $\psi(D_R)+\ell_u=N$.

If $\lambda\le1$, then $N=\psi(Z_R)\le|Z_R|\le\theta$ for any $R$.  Let $\lambda>1$;
a tuple of weight less than $\lambda$ has only letters in $\Sigma_1$.  If
$N<\lambda$, take $u$ the root and $R$ its rule: every $\boldsymbol\eta_{R,j}$ has
weight $\ell_{u_j}<\lambda$, and $R$ writes only letters of $w$, so
$\boldsymbol\eta_R$ is a word of type $1$ derived from the start symbol, hence in
$L(G)$, of length at most $\theta$; thus $N\le\theta$.  If $N\ge\lambda$, choose a
node $u$ with $h:=\ell_u\ge\lambda$ whose children all have $\ell_{u_j}<\lambda$
(descend from the root), and let $R$ be its rule.  Its witness children have only
letters in $\Sigma_1$ and total length $\ell_{u_j}$, and $R$ writes the same
letters in $T$ and in $Z_R$, so $\|\boldsymbol\eta_R\|_1=h$.  Every letter weighs
at most $\lambda$, so $|D_R|_{\mathrm{fix}}\ge\psi(D_R)/\lambda=(N-h)/\lambda$, and
\[
 \theta\ge|Z_R|\ge h+\frac{N-h}{\lambda},\qquad
 N\le\lambda\theta-(\lambda-1)h\le\lambda\theta-(\lambda-1)\lambda
 =\lambda(\theta+1-\lambda)\le\frac{(\theta+1)^2}{4}.
\]
In every case $N\le(\theta+1)^2/4$, and $\theta\le(\theta+1)^2/4$.  Since $N$ is an
integer, $\mu^\star(L(G),s)\le\lfloor(\theta+1)^2/4\rfloor$.
\end{proof}

%==============================================================================
% Appendix K, second half: separating the typing classes by explicit targets (flag separation, LG-needed target, nonabelian-needed target)
%==============================================================================
\subsection{Separating the Typing Classes}
\label{subsec:necessity}

The feature example is covered by a window and an abelian group.  Two further
fan-out-two targets separate the regimes of Section~\ref{sec:char-data}: $L_{\mathrm A}$
is admitted by a local-group typing but by no window--group-determined typing, and
$L_{\mathrm B}$ by a window--$S_3$ typing but by no window--abelian one.  Both have
the form $\varsigma w\mathtt\$x_\times$: a flag $\varsigma=\mathtt r$ leaves $w$ unconstrained,
$\varsigma=\mathtt l$ imposes a typing-dependent condition on $w$, and
$x_\times\in C_\times:=\{\mathtt{1}^m\mathtt{2}^n\mathtt{3}^m\mathtt{4}^n:m,n\ge1\}$ makes the
target non-context-free.  Let $\varphi_\times$ map
$\mathtt 1,\mathtt 3,\mathtt 2,\mathtt 4$ to $(1,0),(-1,0),(0,1),(0,-1)\in\mathbb Z^2$
and all other letters to $0$, and let $\#_{\mathtt{ab}}(w)$ be the number of
occurrences of the factor $\mathtt{ab}$ in $w$.

\begin{lemma}[Flag separation]
\label{lem:flag-separation}
Let $s'$ be a typing, $u,v\in\Sigma^+$ with $s'(u)=s'(v)$, and $z\in\Sigma^*$.  If $L$
contains $\mathtt ruz$, $\mathtt rvz$ and $\mathtt luz$ but not $\mathtt lvz$, then
$L\notin\Cmcf{f}{s'}$ for every $f\ge1$.
\end{lemma}

\begin{proof}
$\mathtt r\blank_1z\in\mathcal E_1$ accepts $u$ and $v$, while $\mathtt l\blank_1z$
accepts $u$ but not $v$; so $\D^{(1)}_L(u)\ne\D^{(1)}_L(v)$, and every $(f,s')$
condition includes arity one.
\end{proof}

\paragraph{The targets.}
Let $\Sigma_{\mathrm A}=\{\mathtt l,\mathtt r,\mathtt a,\mathtt b,\mathtt c,\mathtt\$,
\mathtt 1,\mathtt 2,\mathtt 3,\mathtt 4\}$,
\[
 L_{\mathrm A}:=\{\varsigma w\mathtt\$x_\times:\varsigma\in\{\mathtt l,\mathtt r\},\
 w\in\{\mathtt a,\mathtt b,\mathtt c\}^+,\ x_\times\in C_\times,\
 \varsigma=\mathtt r\text{ or }\#_{\mathtt{ab}}(w)\text{ even}\},
\]
and $s_{\mathrm A}:=\lambda_{1,\nu}|_{\Sigma_{\mathrm A}^+}$ with
$\nu:\Sigma_{\mathrm A}^2\to\mathbb Z_2$, $\nu(\mathtt{ab})=1$ and $\nu=0$ otherwise
(Definition~\ref{def:local-group-typing}); it records the first and last letters and,
for $|w|\ge2$, the parity of $\#_{\mathtt{ab}}(w)$.  Let
$\Sigma_{\mathrm B}=\Sigma_{\mathrm A}\setminus\{\mathtt c\}$, $\alpha=(12)$,
$\beta=(23)\in S_3$, $g_{\mathrm B}:\Sigma_{\mathrm B}^*\to S_3$ with
$\mathtt a\mapsto\alpha$, $\mathtt b\mapsto\beta$ and all other letters $\mapsto1$,
\[
 L_{\mathrm B}:=\{\varsigma w\mathtt\$x_\times:\varsigma\in\{\mathtt l,\mathtt r\},\
 w\in\{\mathtt a,\mathtt b\}^+,\ x_\times\in C_\times,\
 \varsigma=\mathtt r\text{ or }g_{\mathrm B}(w)=1\},
\]
and $s_{\mathrm B}:=(s_{1,1},g_{\mathrm B})|_{\Sigma_{\mathrm B}^+}$, corestricted.
$\mathbf{Ab}$ is the pseudovariety of finite abelian groups, and
$\mathbf{LI}\vee\mathbf{Ab}$ consists of the divisors of $\mathsf T\times {\mathbb G}$ with
$\mathsf T\in\mathbf{LI}$ and ${\mathbb G}$ finite abelian.

\begin{proposition}[A target separating local-group and window--group typings]
\label{prop:lg-needed-target}
\begin{enumerate}[label=\textup{(\roman*)},leftmargin=*]
\item $L_{\mathrm A}\in\mathsf{MCFL}(2,2)\setminus\mathrm{CFL}$.
\item $L_{\mathrm A}\in\Cmcf{f}{s_{\mathrm A}}$ for every $f\ge2$, and
$s_{\mathrm A}(\Sigma_{\mathrm A}^+)\in\mathbf{LG}$.
\item If $s'\preceq(s_{k,\ell},g'|_{\Sigma_{\mathrm A}^+})$ for some $k,\ell$ and a
morphism $g'$ into a finite group (in particular if
$s'(\Sigma_{\mathrm A}^+)\in\mathbf{LI}\vee\mathbf{Grp}$), then
$L_{\mathrm A}\notin\Cmcf{f}{s'}$ for all $f\ge1$.
\end{enumerate}
Hence $s_{\mathrm A}(\Sigma_{\mathrm A}^+)\in\mathbf{LG}\setminus(\mathbf{LI}\vee\mathbf{Grp})$,
and $L_{\mathrm A}$ is covered by Theorem~\ref{thm:local-group-polydata} but, for no
typing, by Theorem~\ref{thm:window-group-polydata}, even with
Lemma~\ref{lem:mcfg-li-group-window}.
\end{proposition}

\begin{proposition}[A target separating nonabelian and window--abelian typings]
\label{prop:nonabelian-needed-target}
\begin{enumerate}[label=\textup{(\roman*)},leftmargin=*]
\item $L_{\mathrm B}\in\mathsf{MCFL}(2,2)\setminus\mathrm{CFL}$.
\item $L_{\mathrm B}\in\Cmcf{f}{s_{\mathrm B}}$ for every $f\ge2$, and $s_{\mathrm B}$ is
window--group-determined with group $S_3$.
\item If $s'\preceq(s_{k,\ell},g'|_{\Sigma_{\mathrm B}^+})$ for some $k,\ell$ and a
morphism $g'$ into a finite \emph{abelian} group (in particular if
$s'(\Sigma_{\mathrm B}^+)\in\mathbf{LI}\vee\mathbf{Ab}$), then
$L_{\mathrm B}\notin\Cmcf{f}{s'}$ for all $f\ge1$.
\end{enumerate}
Hence $L_{\mathrm B}$ is covered by Theorem~\ref{thm:window-group-polydata} with the
nonabelian constant $c_{S_3,f}=M_{S_3,f}-1$, but for no typing by its trivial-group or
abelian-group cases (Remark~\ref{rem:window-abelian}).
\end{proposition}

\begin{proof}[Proof of Propositions~\ref{prop:lg-needed-target}
and~\ref{prop:nonabelian-needed-target}]
Write $\diamond\in\{\mathrm A,\mathrm B\}$, let $\Sigma_\diamond^\circ$ be $\{\mathtt a,\mathtt b,
\mathtt c\}$ or $\{\mathtt a,\mathtt b\}$, and let $P_\diamond(w)$ be ``$\#_{\mathtt{ab}}(w)$
even'' or ``$g_{\mathrm B}(w)=1$''.

(i) Let $W$ be a right-linear (fan-out-one, rank-one, good) grammar for the regular
language $\{\varsigma w:\varsigma\in\{\mathtt l,\mathtt r\},w\in(\Sigma_\diamond^\circ)^+,
\varsigma=\mathtt r\text{ or }P_\diamond(w)\}$; it tracks the last letter and the parity,
respectively the value in $S_3$.  Add
\[
 C\to(\xi_1^1\mathtt\$\mathtt 1,\mathtt 3)(W)\mid(\xi_1^1\mathtt 1,\xi_1^2\mathtt 3)(C),\quad
 D\to(\mathtt 2,\mathtt 4)\mid(\xi_1^1\mathtt 2,\xi_1^2\mathtt 4)(D),\quad
 S\to(\xi_1^1\xi_2^1\xi_1^2\xi_2^2)(C,D).
\]
Then $L_C=\{(\varsigma w\mathtt\$\mathtt 1^m,\mathtt 3^m)\}$, $L_D=\{(\mathtt 2^n,\mathtt 4^n)\}$,
and $S$ generates $L_\diamond$; the grammar is good (adjacent variables in the start
template belong to different children).  Intersecting with
$\mathtt{ra\$}\mathtt 1^+\mathtt 2^+\mathtt 3^+\mathtt 4^+$ and erasing $\mathtt{ra\$}$
gives the non-context-free $C_\times$.

(ii) $L_\diamond=Q_\diamond\cap\Pi_\diamond\cap\varphi_\times^{-1}(0)$ with
$Q_\diamond=\{\mathtt l,\mathtt r\}(\Sigma_\diamond^\circ)^+\mathtt\$\mathtt 1^+\mathtt 2^+\mathtt 3^+
\mathtt 4^+$ and $\Pi_\diamond$ the set of nonempty words that begin with $\mathtt r$, or begin
with $\mathtt l$ and satisfy $P_\diamond$ for the whole word.  On $Q_\diamond$ the condition $P_\diamond$ of
the whole word is that of $w$: no factor $\mathtt{ab}$ meets $\varsigma$, $\mathtt\$$ or
the digits, and $g_{\mathrm B}$ is trivial on all other letters.  $Q_\diamond$ is strictly
2-local (its pairs are $\vdash\varsigma$, $\varsigma x$, $xy$, $x\mathtt\$$, $\mathtt{\$1}$,
$\mathtt{11},\mathtt{12},\mathtt{22},\mathtt{23},\mathtt{33},\mathtt{34},\mathtt{44}$,
$\mathtt 4\dashv$ for $\varsigma\in\{\mathtt l,\mathtt r\}$, $x,y\in\Sigma_\diamond^\circ$) and
$s_{1,1}\preceq s_\diamond$, so Lemma~\ref{lem:sl2-boundary-control} applies.  $\Pi_\diamond$ is a
union of $s_\diamond$-classes: $s_\diamond$ records the first letter, and for words of length at least
two the parity of $\#_{\mathtt{ab}}$, respectively the value of $g_{\mathrm B}$ (the
one-letter words $\mathtt l,\mathtt r$ lie in $\Pi_\diamond$); so item~(c) after
Definition~\ref{def:class-cmcf} applies.  Proposition~\ref{prop:abelian-coset-filter}
handles $\varphi_\times^{-1}(0)$ and the intersection; with (i),
$L_\diamond\in\Cmcf{f}{s_\diamond}$ for $f\ge2$.  $s_{\mathrm A}$ is local-group-determined, so its
positive image lies in $\mathbf{LG}$ by Theorem~\ref{thm:lg-characterization}.

(iii) for $\mathrm A$.  Let $\hat k=\max\{k,\ell,1\}$, $q\ge2$ an even multiple of the
exponent of the group of $g'$, $u=\mathtt a^{\hat k}(\mathtt{abc})^q\mathtt a^{\hat k}$
and $v=\mathtt a^{\hat k}(\mathtt{bac})^q\mathtt a^{\hat k}$.  Both are long for
$s_{k,\ell}$ with prefix $\mathtt a^k$ and suffix $\mathtt a^\ell$, and
$g'(u)=g'(\mathtt a)^{2\hat k}=g'(v)$; so $s'(u)=s'(v)$.  But $\#_{\mathtt{ab}}(u)=q$ is
even and $\#_{\mathtt{ab}}(v)=1$ (only at $\mathtt a^{\hat k}\mathtt b$).  With
$z=\mathtt{\$1234}$, Lemma~\ref{lem:flag-separation} applies.  A typing with positive
image in $\mathbf{LI}\vee\mathbf{Grp}$ satisfies the hypothesis by
Lemma~\ref{lem:mcfg-li-group-window}.

(iii) for $\mathrm B$.  Let $\hat k\ge\max\{k,\ell,1\}$ be even,
$u=\mathtt a^{\hat k}(\mathtt{ab})^3\mathtt a^{\hat k}$ and
$v=\mathtt a^{\hat k}(\mathtt{ab})^2\mathtt{ba}\,\mathtt a^{\hat k}$.  They have the same
prefix $\mathtt a^k$, suffix $\mathtt a^\ell$ and Parikh vector, so $s'(u)=s'(v)$ for abelian
$g'$.  Since $\alpha^2=\beta^2=1$, $\alpha\beta$ has order $3$ and
$\beta\alpha=(\alpha\beta)^{-1}$, we get $g_{\mathrm B}(u)=(\alpha\beta)^3=1$ and
$g_{\mathrm B}(v)=(\alpha\beta)^2(\alpha\beta)^{-1}=\alpha\beta\ne1$.  With
$z=\mathtt{\$1234}$, Lemma~\ref{lem:flag-separation} applies.  For
$\mathbf{LI}\vee\mathbf{Ab}$, the proof of Lemma~\ref{lem:mcfg-li-group-window} with
an abelian group factor gives an abelian $g'$.

\emph{Consequences.}  If $s_{\mathrm A}(\Sigma_{\mathrm A}^+)$ were in
$\mathbf{LI}\vee\mathbf{Grp}$, (iii) with $s'=s_{\mathrm A}$ would contradict (ii).
Theorem~\ref{thm:window-group-polydata} and its trivial-group and abelian cases
assume a typing of the excluded kind.
\end{proof}

%==============================================================================
% Appendix L: rule rank one
%==============================================================================
\section{Rule Rank One}
\label{sec:further}
\label{subsec:rank-one}

At rule rank one, grammar size alone bounds characteristic data under every
typing.  Notation is that of Section~\ref{sec:char-data}; in particular
$C_s(G)=|V||\mathsf S^\sharp|^f+|P||\mathsf S^\sharp|^{fr_{\max}}\le
|\mathsf S^\sharp|^{fr_{\max}}\|G\|$.

\begin{theorem}[Rank-one characteristic-data bound]
\label{thm:rankone-polydata}
Fix $f$ and a typing $s$.
\begin{enumerate}[label=\textup{(\roman*)},leftmargin=*]
\item If $G$ is a reduced good fan-out-$\le f$, rank-$\le1$ presentation of
$L\cap\Sigma^+$ with $L\in\Cmcf{f}{s}$, then
$\Wit(\widetilde G_0)\cup(\{\eps\}\cap L)$ is a characteristic sample for both
$\mathcal B^{\le1}_{f,s}$ and $\mathcal B_{f,s}$, and
\[
 \|\Wit(\widetilde G_0)\|_+\le|\mathsf S^\sharp|^f\bigl(3+|\mathsf S^\sharp|^f\bigr)
 \|G\|^2\bar\theta_G
 \le2|\mathsf S^\sharp|^f\bigl(3+|\mathsf S^\sharp|^f\bigr)\|G\|^3 .
\]
\item For any standard fan-out-$\le f$, rank-$\le1$ grammar $G_{\mathrm{in}}$ for
$L\in\Cmcf{f}{s}$, the operators $\mathcal B^{\le1}_{f,s}$ and $\mathcal B_{f,s}$
have a common characteristic sample $K$ with $\|K\|_+=O_{f,s}(\|G_{\mathrm{in}}\|^3)$.
\end{enumerate}
\end{theorem}

\begin{proof}
(i) A terminal derivation in a rank-$\le1$ grammar is a chain of unary rules
ending in a constant rule.  In a shortest chain no nonterminal repeats (splicing
out the segment between two occurrences does not lengthen the yield, by
linearity), so its rules are distinct.  Hence $t_G\le\|G\|$, and the proof of
Lemma~\ref{lem:thickness-equivalence}(iii), whose path rules together with $R$ are
distinct and whose child chain is distinct, gives $\theta_G\le2\|G\|$, so
$\bar\theta_G\le2\|G\|$.  In $\widetilde G_0$, which also has rank $\le1$, a shortest
chain from a typed nonterminal $X$ repeats no typed nonterminal, so it has at
most $|V||\mathsf S^\sharp|^f$ rules, each writing at most $\bar\theta_G$
letters (Lemma~\ref{lem:thickness-equivalence}(i)).  So
$\|\omega(X)\|_1\le b:=|V||\mathsf S^\sharp|^f\bar\theta_G$ for suitable canonical
tuples.  Lemma~\ref{lem:witness-length-template}(ii),(iv) with $r_{\max}=1$ gives
the sample and
$\|\Wit\|_+\le C_s(G)(1+2\bar\theta_G+b)$; use
$C_s(G)\le|\mathsf S^\sharp|^f\|G\|$, $1+2\bar\theta_G\le3\bar\theta_G$ and
$|V|\le\|G\|$.

(ii) If $L\subseteq\{\eps\}$ take $K=L$.  Otherwise
Proposition~\ref{prop:good-normalization} gives a reduced good rank-$\le1$
presentation $G$ of $L\cap\Sigma^+$ with $\|G\|=O_f(\|G_{\mathrm{in}}\|)$; apply (i).
\end{proof}

Here the special rank-one argument $\bar\theta_G\le2\|G\|$ is
essential for the cubic bound; applying only the general thickness lemma
would lose a power of $\|G\|$.  Equivalently, the shortest typed chains
above give $t_{\mathrm{typ}}\le |V||\mathsf S^\sharp|^f\bar\theta_G$.
No thickness parameter is needed at rank one, and the typing affects only the
constant; this extends Yoshinaka's untyped bound $\|K_G\|\le|P|\theta_G$
\cite[Lemma~7]{Yoshinaka2011}, which at rank one is $O(\|G\|^2)$ since
$\bar\theta_G\le2\|G\|$.  Rank two needs thickness: the size-$O(n)$
grammar $S\to A_n$, $A_i\to A_{i-1}A_{i-1}$, $A_0\to\mathtt c$ generates
$\{\mathtt c^{2^n}\}$, which lies in $\Cmcf{f}{s}$ for every $f$ and $s$ (over one
letter, a shared sentence context forces equal total component length, hence equal
distributions), and every characteristic sample for $\mathcal B^{\le r_{\max}}_{f,s}$
contains $\mathtt c^{2^n}$, because it is a nonempty subset of the target
($L(\mathcal B(\emptyset))=\emptyset$).

%==============================================================================
% Section 5: closing paragraph on target separations
%==============================================================================
\paragraph{The regimes are different on targets.}
Appendix~\ref{sec:comparison} separates the typing regimes by explicit
non-context-free targets: a target that needs the sliding-factor information of
a local group and admits no window--group typing
(Proposition~\ref{prop:lg-needed-target}), and a target that admits a
window--$S_3$ typing but no window--abelian one
(Proposition~\ref{prop:nonabelian-needed-target}).  It also shows that the trivial
typing recovers Yoshinaka's multidimensional substitutability,
$\Cmcf{f}{s_{\mathrm{triv}}}=\YSub{f}\cap f\text{-}\mathrm{MCFL}$.

%==============================================================================
% Appendix K: subsection heading removed when its second half was parked
% (the first half is now the whole appendix, label sec:comparison)
%==============================================================================
\subsection{Yoshinaka's Baseline and Finite Feature Licensing}
\label{subsec:baseline}
